% ----- CHAUPAI 39 -----

\newpage

\subsection*{\deva{एकोनचत्वारिंशत्तम चौपाई} (Thirty-Ninth Chaupai)}
\addcontentsline{toc}{subsection}{\deva{एकोनचत्वारिंशत्तम चौपाई} (Thirty-Ninth Chaupai) :  \deva{जो यह...}}

\begin{minipage}[t]{0.55\textwidth}
\vspace{0pt}

{\large \linespread{1.5}\selectfont
\deva{जो यह पढ़ै हनुमान चलीसा।}\\
\deva{होय सिद्धि साखी गौरीसा॥}
\par}

\vspace{0.5em}

{\large \linespread{1.5}\selectfont
\telu{జో యహ పఢై హనుమాన చలీసా।}\\
\telu{హోయ సిద్ధి సాఖీ గౌరీసా॥}
\par}

\vspace{0.5em}

{\large \linespread{1.3}\selectfont
\textit{jo yaha paṛhai hanumāna calīsā |}\\
\textit{hoya siddhi sākhī gaurīsā ||}
\par}
\end{minipage}%
\hfill
\begin{minipage}[t]{0.42\textwidth}
\vspace{0pt}
\begin{tcolorbox}[anchorbox]
\textbf{The Universal Guarantee:} When you fully commit to the path of discipline, hard work, and good character, the universe itself aligns to ensure your ultimate success (\textit{Hoya Siddhi}). Consistent effort always pays off in the end.
\vspace{0.5em}\newline When you commit to character and hard work, the universe aligns the highest authorities to guarantee your success.\newline\null\hfill\href{https://www.valmiki.iitk.ac.in/sloka?field_kanda_tid=6&language=dv&field_sarga_value=119}{\textit{Yuddha Kanda, 119:1-8}}
\end{tcolorbox}
\end{minipage}

\vspace{0.5em}
\subsubsection*{\textbf{Visual Rhythm Map:}}
\begin{table}[h!]
\centering
\renewcommand{\arraystretch}{1.2}
\noindent\makebox[\textwidth][l]{%
  {%
  \setlength{\tabcolsep}{0pt}%
  \begin{tabularx}{1.125\textwidth}{|*{18}{>{\centering\arraybackslash\hsize=1.0\hsize}X|}}
  \hline
  \multicolumn{2}{|>{\centering\arraybackslash\hsize=2.0\hsize}X|}{\textbf{\guru{जो}} \par (2)} &
  \multicolumn{2}{>{\centering\arraybackslash\hsize=2.0\hsize}X|}{\textbf{\laghu{य}\laghu{ह}} \par (2)} &
  \multicolumn{3}{>{\centering\arraybackslash\hsize=3.0\hsize}X|}{\textbf{\laghu{प}\guru{ढै}} \par (3)} &
  \multicolumn{5}{>{\centering\arraybackslash\hsize=5.0\hsize}X|}{\textbf{\laghu{ह}\laghu{नु}\guru{मा}\laghu{न}} \par (5)} &
  \multicolumn{5}{>{\centering\arraybackslash\hsize=5.0\hsize}X|}{\textbf{\laghu{च}\guru{ली}\guru{सा}} \par (5)} \\
  \hline
  \end{tabularx}%
  }%
  \hspace{0.01\textwidth}%
  \begin{minipage}{0.1\textwidth}
  \vspace{-0.5em}
  \centering
  {\Huge \textbf{17}}
  \end{minipage}%
}
\vspace{-1.5em}
\end{table}

\begin{table}[h!]
\centering
\renewcommand{\arraystretch}{1.2}
\setlength{\tabcolsep}{0pt}
\begin{tabularx}{\textwidth}{|*{16}{>{\centering\arraybackslash\hsize=1.0\hsize}X|}}
\hline
\multicolumn{3}{|>{\centering\arraybackslash\hsize=3.0\hsize}X|}{\textbf{\guru{हो}\laghu{य}} \par (3)} &
\multicolumn{3}{>{\centering\arraybackslash\hsize=3.0\hsize}X|}{\textbf{\guru{सि}\laghu{द्धि}} \par (3)} &
\multicolumn{4}{>{\centering\arraybackslash\hsize=4.0\hsize}X|}{\textbf{\guru{सा}\guru{खी}} \par (4)} &
\multicolumn{6}{>{\centering\arraybackslash\hsize=6.0\hsize}X|}{\textbf{\guru{गौ}\guru{री}\guru{सा}} \par (6)} \\
\hline
\end{tabularx}
\vspace{-1.5em}
\end{table}

\vspace{1em}
\textbf{ \deva{सरल-संस्कृतम्}:}\\
 \deva{यः एतत् "हनुमान-चालीसा" स्तोत्रं पठति}\\
 \deva{सः सिद्धिं प्राप्नोति, अत्र गौरीशः (शिवः) एव साक्षी अस्ति॥}\\


Whoever reads this Hanuman Chalisa will attain perfection and success.
Lord Shiva (husband of Gauri) is witness to this!

\vspace{1em}
\newpage
\textbf{\deva{प्रतिपदार्थः} (Meaning of each word):}
\begin{table}[h!]
\centering
\renewcommand{\arraystretch}{1.2}
\begin{tabularx}{\textwidth}{@{} l l X X @{}}
\toprule
\textbf{Awadhi} & \textbf{Sanskrit} & \textbf{English} & \textbf{Grammar \& Notes} \\
\midrule
\deva{जो यह} / \telu{జో యహ}  &  \deva{यः एतत्}  &  Whoever this  &   \\
\deva{पढ़ै} / \telu{పఢై}  &  \deva{पठति}  &  Reads  & Verb ending in `-ai` \\
\deva{हनुमान चलीसा} / \telu{హనుమాన చలీసా}  &  \deva{हनुमान-चालीसा}  &  Forty verses of Hanuman  &   \\
\deva{होय सिद्धि} / \telu{హోయ సిద్ధి}  &  \deva{सिद्धिः भवति}  &  Success/Perfection happens  &   \\
\deva{साखी} \gotchaicon / \telu{సాఖీ}  &  \deva{साक्षी}  &  Witness  & \\ 
\deva{गौरीसा} / \telu{గౌరీసా}  &  \deva{गौरीशः (शिवः)}  &  Lord of Gauri  & \\ 
\bottomrule
\end{tabularx}
\end{table}

\begin{tcolorbox}[gotchabox]
\begin{itemize}
    \item \textbf{The \deva{क्ष} $\rightarrow$ \deva{ख} Consonant Shift:} The majestic Sanskrit word for Witness is \deva{साक्षी} (Saakshi). Awadhi completely shatters the harsh `kshi` conjunct into the smooth and aspirated \deva{साखी} (Saakhee). 
\end{itemize}
\end{tcolorbox}

\newpage
